\section{Weighted integrability on the unipotent group}
\label{ju-section}

We now pass from the coordinate estimates of
Section~\ref{jk-section} to the unipotent integral that occurs in
the Jacquet majorant.  We prove the required coordinate identities,
including their Jacobians, and the complete induction.
The negative norm of the vector defining the character is used at
a specific point: it prevents the coefficient needed by
Theorem~\ref{jk-oscillatory} from vanishing in the last step.

\subsection{The Hermitian spaces, flags, and unipotent groups}

Fix \(k\geq1\).  Let \(V_0=\mathbb C^{k+1}\), with its usual
Euclidean norm and Hermitian form
\[
 h_0(v,w)=\sum_{j=1}^{k}v_j\overline{w_j}
                         -v_-\overline{w_-}.
\]
All Hermitian forms here are linear in the first variable.
Inductively set
\[
 V_r=\mathbb Ce_0\oplus V_{r-1}\oplus\mathbb Cf_0,\qquad
 h_r(\alpha e_0+v+\beta f_0,\alpha'e_0+w+\beta'f_0)
       =\alpha\overline{\beta'}+\beta\overline{\alpha'}
                                      +h_{r-1}(v,w).
\]
The three summands are Euclidean orthogonal, and \(e_0,f_0\)
have Euclidean norm one.  Relabel the vectors of the smaller
space as \(e_1,\ldots,e_{r-1}\), \(f_1,\ldots,f_{r-1}\).
Thus \(h_r(e_i,f_j)=\delta_{ij}\).
Write \(h_r(v,w)=w^*H_rv\); then \(H_r^*=H_r\) and \(H_r^2=1\).
Let
\[
 G_r=U(V_r,h_r),\qquad K_r=G_r\cap U(V_r)
\]
where the second unitary group uses the Euclidean norm.
Every element of \(K_r\) commutes with \(H_r\).
The subgroup \(N_r\) consists of the elements \(n\in G_r\) with
\[
 ne_j-e_j\in\operatorname{span}(e_0,\ldots,e_{j-1})
                 \quad(0\leq j<r),\qquad
 nw-w\in\operatorname{span}(e_0,\ldots,e_{r-1})
                 \quad(w\in V_0).
\]
In particular \(N_0=\{1\}\).

Let \(u_1\) denote the first positive unit vector in \(V_0\),
and put
\[
 \varepsilon=(u_1+u_-)/\sqrt2,\qquad
 (\phi_0,\ldots,\phi_r)=(f_0,\ldots,f_{r-1},\varepsilon),
 \qquad \phi'_j=H_r\phi_j .
\]
The vectors \(\phi_j,\phi'_j\), together with
\(u_2,\ldots,u_k\), form a Euclidean orthonormal basis.
Their Hermitian products are
\[
 h_r(\phi_i,\phi_j)=h_r(\phi'_i,\phi'_j)=0,\qquad
 h_r(\phi_i,\phi'_j)=\delta_{ij}.
\]
For \(g\in G_r\), define
\[
 \Delta_l(g)=
       \|g\phi_0\wedge\cdots\wedge g\phi_l\|>0
                         \quad(0\leq l\leq r),
 \qquad
 \Xi_s(g)=\prod_{l=0}^{r}\Delta_l(g)^{-s_l}.
\]
The exterior norm is the Euclidean one: its square is the
determinant of the Gram matrix.
The critical exponents are
\[
 s_l^0=\begin{cases}2&l<r,\\k&l=r.\end{cases}
\]

\begin{lemma}\label{ju-heisenberg}
For \(r\geq1\), let \(W=V_{r-1}\) and \(J=H_{r-1}\).
Set
\[
 n(v,t)=
 \begin{pmatrix}
 1&-v^*J&it-\tfrac12 h_{r-1}(v,v)\\
 0&I&v\\0&0&1
 \end{pmatrix},\qquad
 \iota(m)=\operatorname{diag}(1,m,1).
\]
Every element of \(N_r\) has a unique expression
\(n(v,t)\iota(m)\), with \(v\in W\), \(t\in\mathbb R\),
and \(m\in N_{r-1}\).  These are global smooth coordinates.
There is a Haar measure on \(N_r\) which in these coordinates is
\(dv\,dt\,dm\), where \(dm\) is the recursively chosen Haar
measure on \(N_{r-1}\), starting with mass one on \(N_0\).
\end{lemma}

\begin{proof}
The matrix \(n(v,t)\) preserves the Hermitian form by direct
multiplication.  It satisfies the defining flag conditions.
An element \(g\in N_r\) fixes \(e_0\).  Since it preserves the
form, its last row is \((0,0,1)\); consequently
\[
 g=\begin{pmatrix}1&\alpha&\beta\\0&m&v\\0&0&1\end{pmatrix}.
\]
Pairing its middle columns with each other gives \(m^*Jm=J\).
Pairing them with its last column gives
\(\alpha=-v^*Jm\), and the norm of the last column gives
\(2\operatorname{Re}\beta=-h_{r-1}(v,v)\).
The flag conditions on the middle block are exactly those of
\(N_{r-1}\).  Thus \(t=\operatorname{Im}\beta\) gives the
claimed factorization, its uniqueness, and its smooth inverse.

Multiplication of the displayed matrices gives
\[
 n(v,t)n(v',t')=
 n(v+v',\,t+t'-\operatorname{Im}h_{r-1}(v',v)),\qquad
 \iota(m)n(v,t)\iota(m)^{-1}=n(mv,t).
\]
Left and right translations in the first group have real
Jacobian one.  Also \(|\det_{\mathbb C}m|=1\), by taking
determinants in \(m^*Jm=J\), so the real Jacobian of
\((v,t)\mapsto(mv,t)\) is one.
It follows directly from these two formulas and induction that
\(dv\,dt\,dm\) is invariant under both left and right translations.
It is positive on nonempty open sets and finite on compact sets,
so it is a Haar measure.
\end{proof}

\subsection{The exact exterior norm formula}

\begin{lemma}\label{ju-coordinate-formula}
On \(W=V_{r-1}\), write
\[
 v=\sum_{j=0}^{r-1}(a_j\phi'_j+b_j\phi_j)
                        +\sum_{i=2}^{k}z_i u_i,\qquad
 \widetilde t=t+\operatorname{Im}\sum_jb_j\overline{a_j}.
\]
The map \((v,t)\mapsto(a,b,z,\widetilde t)\) preserves
Lebesgue measure.  With \(p=r\), \(m=k-1\), and the
quantities of \eqref{jk-coordinates}, one has
\begin{equation}\label{ju-formula}
 \Delta_l(n(v,t))=d_l(a,b,z,\widetilde t)
                                      \quad(0\leq l\leq r).
\end{equation}
\end{lemma}

\begin{proof}
The change from \(v\) to \((a,b,z)\) is unitary for the
Euclidean norm.  The subsequent change in \(t\) is a translation
depending on those coordinates, and hence has Jacobian one.
As
\[
 h_{r-1}(v,v)=
         2\operatorname{Re}\sum_jb_j\overline{a_j}+\|z\|^2,
\]
the upper right entry of \(n(v,t)\) is precisely \(B_0\).
Its relevant column vectors are
\[
 n(v,t)f_0=B_0e_0+v+f_0,\qquad
 n(v,t)\phi_j=\phi_j-\overline{a_j}e_0.
\]
Here the \(\phi_j\) on the right belong to \(W\), and correspond
to indices \(j+1\) in \(V_r\).
In the wedge with \(l\) such middle columns, subtract
\(\sum_{j<l}b_jn(v,t)\phi_j\) from the first column.
The first column becomes
\[
 y=B_le_0+f_0+\sum_j a_j\phi'_j+
                   \sum_{j\geq l}b_j\phi_j+\sum_i z_i u_i.
\]
The remaining columns are \(u_j=\phi_j-\overline{a_j}e_0\),
\(j<l\).
Let \(a_{<l}=(a_0,\ldots,a_{l-1})^T\).
Their Gram matrix, with \(y\) first, is
\[
 \begin{pmatrix}
 R_l+|B_l|^2&-\overline{B_l}\,a_{<l}^*\\
 -B_l a_{<l}&I+a_{<l}a_{<l}^*
 \end{pmatrix}.
\]
The lower block has determinant \(L_l\) and inverse
\(I-a_{<l}a_{<l}^*/L_l\).
The inverse formula follows by multiplication; the determinant
follows by using a unitary basis whose first vector is proportional
to \(a_{<l}\), or immediately when this vector is zero.
Taking the Schur complement now gives
\[
 L_l\left(R_l+|B_l|^2
       -|B_l|^2\frac{\|a_{<l}\|^2}{L_l}\right)
                                  =L_lR_l+|B_l|^2.
\]
For \(l=0\) the same identity uses the empty lower block.
This proves \eqref{ju-formula}.
\end{proof}

\subsection{Orthogonalization of an isotropic flag}

\begin{lemma}\label{ju-isotropic-qr}
Let \(Y_0,\ldots,Y_r\in V_r\) be linearly independent and
span a totally isotropic subspace.  There exist
\(\kappa\in K_r\) and an upper triangular matrix
\(\tau=(\tau_{ij})_{0\leq i,j\leq r}\), with positive real
diagonal entries, such that
\[
 Y_j=\sum_{i\leq j}\tau_{ij}\kappa\phi_i,\qquad
 \|Y_0\wedge\cdots\wedge Y_l\|
                            =\prod_{j=0}^{l}\tau_{jj}.
\]
\end{lemma}

\begin{proof}
Apply Euclidean Gram--Schmidt to the \(Y_j\).
It gives Euclidean orthonormal vectors \(q_j\) in their span
and the indicated upper triangular coefficients with positive
diagonal.  Their span is still totally isotropic.
Let \(V_r=V_+\oplus V_-\) be the orthogonal decomposition
into the eigenspaces of \(H_r\), and let \(P_\pm\) be the
corresponding projections.  The two identities
\[
 \langle q_i,q_j\rangle=\delta_{ij},\qquad h_r(q_j,q_i)=0
\]
imply
\[
 \langle P_+q_i,P_+q_j\rangle
   =\langle P_-q_i,P_-q_j\rangle=\delta_{ij}/2.
\]
The same is true for the \(\phi_j\).
There are \(r+1=\dim V_-\) negative projected vectors,
and \(\dim V_+=k+r\geq r+1\).
Extend the positive projected frames to orthonormal bases;
the negative projected frames are already bases.
The two resulting unitary changes of basis give
\(\kappa=\kappa_+\oplus\kappa_-\in K_r\) with
\(\kappa\phi_j=q_j\).
Taking the wedge of the first \(l+1\) triangular expansions
proves the norm formula.
Since \(H_r\in K_r\) interchanges \(\phi_j\) and \(\phi'_j\),
replacing \(\kappa\) by \(\kappa H_r\) gives the identical
assertion with the primed reference frame.
\end{proof}

\begin{lemma}\label{ju-levi-factor}
For \(m\in G_{r-1}\), one can choose
\(\kappa\in K_{r-1}\) such that
\[
 m\phi_0=\Delta_0(m)\kappa\phi_0 .
\]
For that choice, putting \(v'=\kappa^{-1}v\), one has
\begin{align}
 \Delta_0(n(v,t)\iota(m))&=\Delta_0(n(v',t)),\nonumber\\
 \Delta_{j+1}(n(v,t)\iota(m))
     &=\Delta_j(m)\Delta_{j+1}(n(v',t))
                                    \quad(0\leq j<r).
 \label{ju-levi-norms}
\end{align}
\end{lemma}

\begin{proof}
Apply Lemma~\ref{ju-isotropic-qr} on \(W\) to \(m\phi_j\).
Conjugation gives
\(n(v,t)\iota(\kappa)=\iota(\kappa)n(\kappa^{-1}v,t)\).
The first column \(f_0\) is fixed by \(\iota(m)\).
In the first \(j+2\) columns of the wedge, the triangular
expansion of the other \(j+1\) columns therefore contributes
the factor \(\prod_{i\leq j}\tau_{ii}=\Delta_j(m)\).
The exterior action of \(\iota(\kappa)\) is Euclidean unitary,
so it does not change the norm.
The assertion for the first column follows by the same
unitary invariance.  Finally
\(\tau_{00}=\|m\phi_0\|=\Delta_0(m)\).
\end{proof}

\subsection{The character and the nonvanishing coefficient}

Fix a vector \(z\in V_0\) with \(h_0(z,z)<0\).
For \(r\geq1\), put
\[
 \ell_z(g)=\sum_{a=0}^{r-2}(ge_{a+1})_{e_a}
                                  +(gz)_{e_{r-1}},
 \qquad \lambda_z(g)=2\operatorname{Im}\ell_z(g).
\]
For \(r=0\), put \(\ell_z=\lambda_z=0\).
This is the character normalization in the source.
The flag conditions show that each consecutive entry, and the
last coefficient involving \(z\), is additive under multiplication
in \(N_r\): in multiplying consecutive entries, all possible
intermediate terms lie strictly earlier in the flag.
For the last coefficient use
\(n'z-z\in\operatorname{span}(e_0,\ldots,e_{r-1})\)
and the fact that the coefficient of \(e_{r-1}\) in \(ne_j\)
is \(\mathbf1_{\{j=r-1\}}\).
Thus \(\lambda_z\) is additive on \(N_r\), and
\(\exp(i\lambda_z)\) is a unitary character.
For \(W=V_{r-1}\), let \(f_z=e_0\) if \(r-1\geq1\),
and \(f_z=z\) if \(r-1=0\).
Inspection of the first row and the middle block gives
\begin{equation}\label{ju-lambda}
 \lambda_z(n(v,t)\iota(m))
       =2\operatorname{Im}h_{r-1}(v,mf_z)+\lambda_z(m).
\end{equation}
Indeed the first summand of \(\ell_z\) is
\(-h_{r-1}(mf_z,v)\), and its remaining summands are
exactly \(\ell_z(m)\).

The fixed number
\[
 c_z=h_{r-1}(\phi_0,f_z)
\]
is nonzero.  If \(r-1\geq1\), it equals
\(h_{r-1}(f_0,e_0)=1\).
If \(r-1=0\), it equals \(h_0(\varepsilon,z)\).
Were this zero, the first positive coordinate of \(z\)
would equal its negative coordinate; consequently
\(h_0(z,z)=\sum_{i=2}^{k}|z_i|^2\geq0\), a contradiction.

\begin{proposition}\label{ju-inner}
Let \(0<\delta\leq1/2\), \(0<\gamma\leq\delta/8\),
and \(s_l=s_l^0-e_l\), where \(e_l>0\) and
\(e_0+2\sum_{l=1}^{r}e_l<\gamma/2\).
Define
\[
 s'_j=s_{j+1}-\gamma\,\mathbf1_{\{j=0\}}
                                      \quad(0\leq j<r).
\]
There is a constant independent of \(m\in G_{r-1}\) such that
\begin{align}
 &\int_{W\times\mathbb R}
  \Xi_s(n(v,t)\iota(m))
  (1+|\lambda_z(n(v,t)\iota(m))|)^{-\delta}\,dv\,dt
  \nonumber\\
 &\hspace{15mm}\leq
 C|c_z|^{-\gamma}\Xi_{s'}(m)
                            (1+|\lambda_z(m)|)^{-\delta/4}.
 \label{ju-inner-bound}
\end{align}
\end{proposition}

\begin{proof}
Choose \(\kappa\) as in Lemma~\ref{ju-levi-factor}, set
\(v'=\kappa^{-1}v\), and put \(w=\kappa^{-1}mf_z\).
This change of variables preserves Lebesgue measure.
By \eqref{ju-levi-norms}, the factor independent of \(v',t\)
is \(\prod_{j=0}^{r-1}\Delta_j(m)^{-s_{j+1}}\).
Lemma~\ref{ju-coordinate-formula} gives the remaining
coordinate kernel of Theorem~\ref{jk-oscillatory}, with
\(P=r\), \(m\) there equal to \(k-1\).
In those coordinates,
\[
 h_{r-1}(v',w)=
 \sum_j a_j h_{r-1}(\phi'_j,w)
 +\sum_j b_j h_{r-1}(\phi_j,w)
 +\sum_i z_i h_{r-1}(u_i,w).
\]
Thus the coefficient denoted by \(\eta'_0\) in that theorem
is \(h_{r-1}(\phi_0,w)\).
Since \(m\phi_0=\Delta_0(m)\kappa\phi_0\), invariance of the
Hermitian form gives
\[
 \Delta_0(m)h_{r-1}(\phi_0,w)
            =h_{r-1}(m\phi_0,mf_z)=c_z .
\]
This coefficient is nonzero, and its negative \(\gamma\)-power
is \(|c_z|^{-\gamma}\Delta_0(m)^\gamma\).
Apply Theorem~\ref{jk-oscillatory} with \(c_0=\lambda_z(m)\)
and use \eqref{ju-lambda}.  Combining the powers of
\(\Delta_j(m)\) gives exactly \(\Xi_{s'}(m)\).
No measurable choice of \(\kappa\) as a function of \(m\)
is needed: the argument gives a bound for each fixed \(m\)
by the displayed continuous function of \(m\).
\end{proof}

\begin{theorem}\label{ju-integrability}
For every \(r\geq0\) and \(\delta>0\), there is
\(\epsilon_r(\delta)>0\) such that
\begin{equation}\label{ju-final}
 \int_{N_r}\Xi_s(n)(1+|\lambda_z(n)|)^{-\delta}\,dn<\infty
\end{equation}
whenever \(0<s_l^0-s_l\leq\epsilon_r(\delta)\) for all \(l\).
The same exponent window works for every negative norm vector
\(z\in V_0\).  The integral is finite for every Haar measure on
\(N_r\).
\end{theorem}

\begin{proof}
It suffices to prove the assertion for \(0<\delta\leq1/2\),
because increasing \(\delta\) decreases the integrand.
Take \(\epsilon_0(\delta)=1/10\); the base group is a singleton.
Suppose the result is known at rank \(r-1\) for exponent
\(\delta/4\), and put
\[
 E=\epsilon_{r-1}(\delta/4),\qquad
 \gamma=\min(\delta/8,E/2),\qquad
 \epsilon_r(\delta)=\min\bigl(E/2,\gamma/(4r+4)\bigr).
\]
If \(0<e_l\leq\epsilon_r(\delta)\), then
\[
 e_0+2\sum_{l=1}^{r}e_l
 \leq(2r+1)\frac{\gamma}{4r+4}<\gamma/2.
\]
The lower exponents in Proposition~\ref{ju-inner} have losses
\[
 s_j^0-s'_j=e_{j+1}+\gamma\mathbf1_{\{j=0\}},
\]
where \(s_j^0\) now refers to rank \(r-1\).
These losses are positive and at most \(E\).
Use the Haar coordinates of Lemma~\ref{ju-heisenberg},
Tonelli, and \eqref{ju-inner-bound}.  The induction hypothesis
makes the resulting integral over \(N_{r-1}\) finite.
The constants involving \(z\) are finite since \(c_z\ne0\);
the choice of \(\epsilon_r(\delta)\) does not depend on \(z\).
Finally every Haar measure is a positive scalar multiple of
the one constructed in Lemma~\ref{ju-heisenberg}, by uniqueness
of Haar measure, so finiteness is independent of normalization.
\end{proof}

\subsection{The exact scope of the conclusion}

Theorem~\ref{ju-integrability} establishes the unipotent
integrability estimate of~\cite{alpo-uniform-source} with the
stated flag, character, Hermitian form and critical exponents.
Its proof supplies the required formula and compact
orthogonalization.  The Gram determinant calculation gives a
short direct replacement for the coordinate computation.
The induction retains precisely the exponent change
\(\delta\mapsto\delta/4\) and the subtraction of \(\gamma\)
from the first lower exponent.

The corresponding group integral, the separated
matrix coefficient estimate, and the final dominated
unfolding argument for the Jacquet pairing are additional
steps.  They are not asserted by \eqref{ju-final} alone.


